\newpage
\section{Estructura de los Métodos de Búsqueda}
\label{sec:algoritmos}

Una vez definidos los operadores comunes y la representación del problema, en esta sección se detalla la estructura algorítmica específica de las tres metaheurísticas implementadas: Búsqueda Aleatoria, Búsqueda Voraz (Greedy) y Búsqueda Local.

\subsection{Búsqueda Aleatoria (Random Search)}
La Búsqueda Aleatoria actúa como línea base (\textit{baseline}) estocástica. Su comportamiento consiste en generar iterativamente soluciones al azar dentro del espacio factible, reparándolas y evaluándolas, conservando únicamente la mejor solución encontrada tras agotar el presupuesto de evaluaciones.

El proceso de generación de una solución aleatoria implica asignar a cada peso un valor uniformemente distribuido entre $lo$ y $hi$, inyectar un número aleatorio de ceros (para garantizar la diversificación exigida por el problema) y, finalmente, aplicar el operador de reparación \texttt{fix(w)}.

\begin{algorithm}[H]
\SetAlgoLined
\KwIn{Problema $P$, máximo de evaluaciones $max\_evals$}
\KwOut{Mejor solución encontrada $w_{best}$ y su $fitness$}
 $w_{best} \leftarrow \emptyset$\;
 $f_{best} \leftarrow -\infty$\;
 \For{$i \leftarrow 1$ \KwTo $max\_evals$}{
  $w_{actual} \leftarrow Generar\_Solucion\_Aleatoria(P)$\;
  \tcp{La generación ya incluye la llamada a fix($w_{actual}$)}
  $f_{actual} \leftarrow fitness(w_{actual})$\;
  \If{$i == 1$ \textbf{or} $f_{actual} > f_{best}$}{
   $w_{best} \leftarrow w_{actual}$\;
   $f_{best} \leftarrow f_{actual}$\;
  }
 }
 \Return $(w_{best}, f_{best})$\;
 \caption{Algoritmo de Búsqueda Aleatoria (Random Search)}
\end{algorithm}

\subsection{Búsqueda Voraz (Greedy Search)}
A diferencia del método anterior, la Búsqueda Voraz es un algoritmo constructivo y determinista que no emplea el azar. Utiliza el conocimiento del problema (heurística) para tomar decisiones localmente óptimas en la construcción de la cartera.

La función heurística evalúa la rentabilidad individual de un activo penalizada por su covarianza media respecto al resto del mercado. Se calcula como:
\begin{equation}
    heur(i) = \left( \sum_{t} \log(1+\mu_{it}) \right) - \lambda \cdot \frac{\sum_{t}\Sigma_{ti}}{T}
\end{equation}

\begin{algorithm}[H]
\SetAlgoLined
\KwIn{Problema $P$, presupuesto $1.0$}
\KwOut{Solución construida $w_{best}$ y su $fitness$}
 $w_{best} \leftarrow \{0.0, \dots, 0.0\}$\;
 $suma\_w \leftarrow 0.0$\;
 Calcular $heur(i)$ para cada empresa $i \in \{1..N\}$\;
 $Ranking \leftarrow$ Ordenar empresas descendentemente según $heur(i)$\;
 \For{cada empresa $c$ en $Ranking$}{
  \If{$suma\_w < 1.0$}{
   $asignar \leftarrow \min(hi, 1.0 - suma\_w)$\;
   $w_{best}[c] \leftarrow asignar$\;
   $suma\_w \leftarrow suma\_w + asignar$\;
  }
  \Else{ \textbf{break}\; }
 }
 \tcp{Garantizar factibilidad de remanentes (ej. $\min < lo$)}
 \texttt{fix}($w_{best}$)\;
 \Return $(w_{best}, fitness(w_{best}))$\;
 \caption{Algoritmo Constructivo Greedy}
\end{algorithm}
\vspace{1cm}
En relación a la factibilidad de las soluciones construidas por este método voraz, cabe destacar que, dadas las características de las instancias proporcionadas (donde el presupuesto de 1.0 es perfectamente divisible por los límites $hi$ establecidos en el IBEX 35, S\&P 100 y S\&P 500, o bien, lo restante no viola las restricciones), el algoritmo agota el capital exacto sin generar remanentes problemáticos. No obstante, para garantizar la escalabilidad y robustez del código ante cualquier configuración de mercado futuro, se ha invocado el operador general de reparación $fix(w)$ al finalizar la construcción, asegurando matemáticamente la viabilidad de la cartera.\\

Por ejemplo, considérese un escenario donde $1.0 / hi = 12.5$ empresas caben al máximo: las primeras 12 empresas ocupan $12 \times 0.083 = 0.996$, dejando un remanente de $0.004$. Si este remanente se asigna directamente a una empresa sin invocar $fix(w)$, violaría la restricción $0.004 < lo = 0.005$, causando un error. La función $fix(w)$ redistribuye automáticamente los pesos para garantizar que todas las asignaciones (tanto las no nulas como el remanente) respetan los límites $[lo, hi]$, asegurando la validez de cualquier solución independientemente de la configuración del mercado. Comentando el fix o no, el resultado es el mismo, lo dejo para cubrir estos casos y reflejar el estudio del algoritmo. \\

\textit{Nota de diseño sobre los parámetros del algoritmo:}
En la descripción formal del algoritmo (Greedy), se ha omitido intencionadamente el parámetro de ``evaluaciones máximas'' (\texttt{max\_evals}) que sí está presente en la implementación del código fuente en C++. Esto se debe a la naturaleza determinista y constructiva del método. A diferencia de las búsquedas estocásticas o basadas en trayectorias, el algoritmo Greedy no explora el espacio de búsqueda de forma iterativa, sino que construye una única solución paso a paso utilizando el conocimiento del problema mediante su función heurística. En consecuencia, el concepto de ``presupuesto de evaluaciones'' carece de sentido teórico, y su existencia en el código responde exclusivamente a la necesidad de mantener el polimorfismo al heredar de la interfaz abstracta común de metaheurísticas. Por el contrario, se ha explicitado el parámetro ``presupuesto'' ($1.0$) en el pseudocódigo para representar fielmente la restricción matemática fundamental del problema: la suma de los pesos de la cartera debe alcanzar estrictamente la unidad.

\subsection{Búsqueda Local (Local Search)}
La Búsqueda Local es un método basado en trayectorias. A partir de una solución inicial generada aleatoriamente, el algoritmo explora su vecindario aplicando un operador de movimiento. Se ha optado por una estrategia de "Primer Mejor" (\textit{First Improvement}), donde el vecindario se baraja aleatoriamente y se acepta el primer vecino que mejore estrictamente el \textit{fitness} actual\footnote{Para ello me he apoyado en la teoría vista tanto en clase de teoría como de prácticas, realizando las modificaciones necesarias.}.

El operador de generación de vecino consiste en seleccionar un par de empresas $(i, j)$, detrayendo un $40\%$ de la inversión actual de la empresa donante $i$ para transferirlo a la receptora $j$. Este movimiento preserva el sumatorio del presupuesto (siempre suma 1), por lo que solo es necesario validar matemáticamente que los nuevos pesos no excedan los límites $[lo, hi]$.

\begin{algorithm}[H]
\SetAlgoLined
\KwIn{Problema $P$, máximo de evaluaciones $max\_evals$}
\KwOut{Óptimo local encontrado $w_{actual}$ y su $fitness$}
 $w_{actual} \leftarrow Generar\_Solucion\_Aleatoria(P)$\;
 $f_{actual} \leftarrow fitness(w_{actual})$\;
 $evals \leftarrow 1$\;
 $Entorno \leftarrow$ Generar todos los pares $(i, j)$ con $i \neq j$\;
 $mejora \leftarrow true$\;
 \While{$evals < max\_evals$ \textbf{and} $mejora == true$}{
  $mejora \leftarrow false$\;
  BarajarAleatoriamente($Entorno$)\;
  \For{cada par $(i, j)$ en $Entorno$}{
   \If{$evals \geq max\_evals$}{ \textbf{break}\; }
   \If{$w_{actual}[i] == 0.0$}{ \textbf{continue}\; }
   
   $trasvase \leftarrow w_{actual}[i] \cdot 0.4$\;
   $nuevo\_i \leftarrow w_{actual}[i] - trasvase$\;
   $nuevo\_j \leftarrow w_{actual}[j] + trasvase$\;
   
   \If{$(nuevo\_i \in [lo, hi] \lor nuevo\_i \approx 0)$ \textbf{and} $(nuevo\_j \in [lo, hi] \lor nuevo\_j \approx 0)$}{
    $w_{vecino} \leftarrow w_{actual}$\;
    $w_{vecino}[i] \leftarrow nuevo\_i$\;
    $w_{vecino}[j] \leftarrow nuevo\_j$\;
    $f_{vecino} \leftarrow fitness(w_{vecino})$\;
    $evals \leftarrow evals + 1$\;
    
    \If{$f_{vecino} > f_{actual}$}{
     $w_{actual} \leftarrow w_{vecino}$\;
     $f_{actual} \leftarrow f_{vecino}$\;
     $mejora \leftarrow true$\;
     \textbf{break}\; \tcp{Primer Mejor encontrado}
    }
   }
  }
 }
 \Return $(w_{actual}, f_{actual})$\;
 \caption{Algoritmo de Búsqueda Local (First Improvement)}
\end{algorithm}